% ============================================================================
\section{Discussion}
\label{sec:discussion}
% ============================================================================
% FIRST DRAFT. Every number from this project is printed by scripts 04-08
% (VERIFICATION.txt). Every statement about a paper was checked against the
% paper. Derived numbers: 247 and 181 CX acting on the average qubit are twice
% Farrell et al.'s CNOT counts (13,858 and 10,154) over 112 qubits; their Table
% III caption prints the same numbers as "twice the number of CNOTs/qubit",
% because they count each CNOT once per circuit, not once per qubit it touches;
% 288 per qubit is 2 x 576 / 4; Shaw et al.'s worst-case CX error is the
% non-identity Pauli probability of the depolarising channel, 15/16 x 4/3 eps2
% = 1.25 eps2, so their bound fails at eps2 n_CX = 0.5 / 1.25 = 0.4.

\subsection{The depth rule and published experiments}

Each of the four qubits takes part in $n_{\mathrm{CX}}/2$ CX gates on average,
so the 1000-shot rule $\epsilon^\ast n_{\mathrm{CX}} \approx 2$ also reads per
qubit: without mitigation the signal is lost once the summed infidelity of the
gates on the average qubit reaches about one.

The per-qubit form is the one to carry to larger devices, since Farrell
\emph{et al.}\ argue that two-qubit gates per qubit, not their total, limit
large simulations~\cite{farrell2024}. Their deepest hadron circuit has
\num{13858} CX gates on 112 qubits, so 247 act on the average
qubit~\cite{farrell2024b}, close to the 288 of the $\delta t = 0.1$ trough
circuit here. With 181 on the average qubit their data showed no wave packet
until heavy error mitigation restored it, and the results still often differed from classical
simulations by many standard deviations~\cite{farrell2024b}. Neither Farrell
paper quotes its device's gate error, and Martinez \emph{et al.}\ give no gate
fidelity and use no CX gates~\cite{martinez2016}, so none of the three can test
the rule. Rosanowski \emph{et al.}\ quote a typical two-qubit fidelity of
99.6(2)\,\% for their trapped-ion processor~\cite{rosanowski2025}; its error,
$2\times10^{-3}$ to $6\times10^{-3}$, straddles $\epsilon^\ast$ at
$\delta t = 0.1$ and lies below it at 0.2 and 0.3.

Shaw \emph{et al.}\ bound the gate-error bias by the number of noisy CX gates
times their worst-case error, $1.25\,\epsilon_2$ for depolarising
CX~\cite{shaw2020}. On the positron density, which cannot exceed $\tfrac12$,
the bound says nothing beyond $\epsilon_2 n_{\mathrm{CX}} = 0.4$, a fifth of
where the dip here is lost: the small-scale comparison the authors ask for,
though their lattice keeps the gauge field.

\subsection{The optimal step}

Knee and Munro~\cite{knee2015} and Shaw \emph{et al.}~\cite{shaw2020} predict
that the best Trotter step grows with the gate error, and the optima here agree:
five, four and three steps to $t_{\mathrm{t}}$ for $\epsilon_2 =
3\times10^{-4}$, $10^{-3}$ and $3\times10^{-3}$, too few integer points to test
a power law. Experiments have made the same choice by eye: Klco \emph{et al.}\
moved from $\delta t = 0.1$ to $0.25$ and Nguyen \emph{et al.}\ doubled their
step, each because after a few steps hardware error, not Trotter error,
dominated~\cite{klco2018,nguyen2022}. At $\epsilon_2 = 3\times10^{-4}$ the
optimal step, 0.35, sits at the 1000-shot crossover at $t_{\mathrm{t}}$, 0.36;
at larger error it is coarser, so at the best achievable step the Trotter error
exceeds the shot noise.

\subsection{What a symmetry check can and cannot tell}

Martinez \emph{et al.}\ called the leakage out of the neutral states a useful
measure of fidelity, and Nguyen \emph{et al.}\ reported it, but neither turned a
charge violation into an error on an observable~\cite{martinez2016,nguyen2022}.
Section~\ref{sec:results} does.

The prediction is optimistic because Eq.~\eqref{eq:FQ} assumes a global
depolarising channel, as do the rescaling methods of Vovrosh \emph{et al.}\ and
Urbanek \emph{et al.}, who treat it as an approximation for deep circuits on
many qubits~\cite{vovrosh2021,urbanek2021}. Four qubits are far from that
regime, and the optimistic bias is the residual: local noise erodes the dip
faster than a global channel with the same charge violation. Farrell \emph{et
al.}\ replaced a rescaling that assumes every qubit decoheres at the same rate
by one factor per observable~\cite{farrell2024}, and still attribute part of
their disagreement with simulation to the simplicity of the assumed
noise~\cite{farrell2024b}. Since every prediction from unsampled
data overshot the true threshold, a threshold read from a symmetry check is best
treated as an upper bound on the usable depth, though the prediction from a
single 1000-shot run falls below the true threshold 8 to 20\,\% of the time.

The blind spot is general, since a symmetry check cannot catch an error that
commutes with the symmetry~\cite{cai2023}, and the noise model of the first
experiment falls in it. Martinez \emph{et al.}\ fitted a phase-flip probability
per qubit per step to their postselected data, choosing that model because
postselection already partly corrects population errors~\cite{martinez2016}.
Their leakage of 14 to 31\,\% and their fitted dephasing thus describe
complementary parts of the noise, and $\mathrm{Var}(Q)$ would read exactly zero
for the second.

\subsection{The ground state}

Rosanowski \emph{et al.}\ argue that SKQD resists noise because only which
bitstrings are sampled matters, provided every needed one appears at least
once~\cite{rosanowski2025}. The exact calculation confirms it: gate error at
$\epsilon_2 = 3\times10^{-3}$ costs only 15 shots per step instead of 12. The
neutral sector grows as $\binom{N}{N/2}$, so rare bitstrings left
unsampled~\cite{miura2026} will limit larger lattices.

\subsection{Limitations}
\label{sec:limitations}

The noise models are the main limitation. Depolarising and dephasing noise act
on each gate independently and without memory; real devices also have coherent
errors such as over-rotation, which could shift or imitate the dip, correlated
errors, and amplitude damping, whose end state Eq.~\eqref{eq:FQ} reads as
perfect fidelity (Section~\ref{sec:noise}). Farrell \emph{et al.}\ saw
systematic errors that grow with the rotation angle of the
gates~\cite{farrell2024b}, which no depolarising model produces and which would
favour smaller steps. Readout error is left out of the threshold; all qubit
pairs are coupled, whereas nearest-neighbour hardware needs SWAP gates for the
RZZ between sites 1 and 3; and the criterion $D \le 2\sigma_D$ is a choice.
Everything is at one parameter point on four sites (twelve for the classical
check of Figure~\ref{fig:sizes}), a coarse lattice far from the continuum, in
simulation with assumed noise.

% ============================================================================
\section{Conclusion}
\label{sec:conclusion}
% ============================================================================

The four-site lattice Schwinger model was simulated with Trotter error, shot
noise, gate error and readout error separated, every stage checked against
exact diagonalisation, and the published curves of Martinez \emph{et al.}\
reproduced. At 1000 shots the Trotter error on the particle density falls
below the shot noise at a step of 0.18 at $t = 3$. At the same shot count the
pair-production signal is lost when the two-qubit infidelity times the CX count
of the deepest circuit reaches about two, at all three step sizes tried, and the
best step grows with the gate error. The charge variance predicts that threshold
with an optimistic bias of 10 to 17\,\% for depolarising noise and a larger one
for dephasing on the gates, and misses it entirely for dephasing between time
steps, the model fitted to the first experiment on this system. Whether these
numbers survive coherent and correlated errors, and real hardware, is the next
question.
